\documentclass[11pt]{article}
\usepackage[utf8]{inputenc}
\usepackage[margin=1in]{geometry}
\usepackage{xcolor}
\usepackage{listings}
\usepackage{hyperref}
\usepackage{titlesec}
\usepackage{array}
\usepackage{enumitem}
\usepackage{amssymb}
\usepackage{listings}
\usepackage{xcolor}

% Configuration for code snippets (shell/bash style)
\lstset{
    backgroundcolor=\color{gray!10},
    basicstyle=\ttfamily\small,
    breaklines=true,
    frame=single,
    rulecolor=\color{black!30},
    keywordstyle=\color{blue},
    commentstyle=\color{green!50!black}
}

% Hyperlink styling
\hypersetup{
    colorlinks=true,
    linkcolor=blue,
    urlcolor=blue,
    citecolor=blue
}

% Title Formatting
\titleformat{\section}{\normalfont\Large\bfseries}{}{0pt}{}
\titleformat{\subsection}{\normalfont\large\bfseries}{}{0pt}{}

% NOTE TO COURSE STAFF: Replace the Ed Discussion thread name in the
% Submission section with the actual thread you want students posting to
% before publishing.

\begin{document}

% --- HEADER SECTION ---
\begin{center}
    {\Huge \textbf{CS193}} \\
    \vspace{2mm}
    {\Large \textbf{Homework - Week 7}} \\
\end{center}

% --- INTRODUCTION ---
\section{Introduction}
In lecture you learned how to go beyond typing one command at a time and start \textbf{scripting} in bash. This week you'll put those skills to work by writing a small but genuinely useful tool: a \textbf{study-log helper}. Each time you run it, it records a timestamped study session for you and reports how many sessions you've logged so far.
\\\\
You only need \textbf{bash and a terminal} --- no other languages, no installs.
\\\\
We are giving you a working starter template. You are not writing the whole script from scratch. The scaffold is laid out with \texttt{TODO} comments, and each one points you to the exact lecture slide that shows the syntax you need. The goal is to leave with a script that runs and that cements variables, arguments, test operators, redirection, and command substitution. It should take about \textbf{15--20 minutes}.

% --- GETTING STARTED ---
\section{Getting Started: The Template}
Open the \texttt{hw7-studylog} folder. It already has the structure and hints you need --- you just fill in the \texttt{TODO} sections.
\begin{itemize}
    \item \texttt{studylog.sh} --- the script you'll write. Look for \texttt{\# TODO} comments. Each one tells you what to add and names the lecture slide with the matching syntax.
    \item \texttt{README.md} --- a quick reminder of how to make the script executable and run it.
    \item \texttt{.vscode/} --- recommended extensions and settings. You don't need to touch this; it just makes editing shell scripts nicer.
    \item \texttt{study\_log.txt} --- you won't see this at first. Your script \textit{creates} it the first time you run it.
\end{itemize}

% --- HOW TO RUN ---
\section{How to Run Your Script}
\subsection{Using VS Code}
\begin{enumerate}
    \item Open VS Code and go to \textbf{File \(\rightarrow\) Open Folder}, then select the \texttt{hw7-studylog} folder.
    \item Open a terminal inside VS Code with \textbf{Terminal \(\rightarrow\) New Terminal}. It opens already pointed at your project folder.
    \item Make the script executable (one time), then run it, passing your name:
\begin{lstlisting}[language=bash]
$ chmod +x studylog.sh
$ ./studylog.sh Boiler
Thanks Boiler! You have logged 1 study session(s).
\end{lstlisting}
    \item Run it a few more times and watch the count grow. Open \texttt{study\_log.txt} to see your entries.
    \item Not using VS Code? The same idea applies in any terminal: \texttt{cd} into the \texttt{hw7-studylog} folder, then run the two commands above. No server, compiler, or installation is required --- it's plain bash.
\end{enumerate}

% --- WHAT TO FILL IN ---
\section{What to Fill In}
Work through the \texttt{TODO} comments in \texttt{studylog.sh} in order:
\begin{enumerate}
    \item \textbf{Shebang.} Keep \texttt{\#!/bin/bash} as the very first line (already done for you). It tells the OS which interpreter to use.
    \item \textbf{Check for a name.} If the student runs the script with no argument, print a usage message and \textbf{exit with a non-zero code}. Hint: test with \texttt{-z "\$1"} inside an \texttt{if}.
    \item \textbf{Use a variable.} Store the log file name in a variable, e.g. \texttt{LOGFILE="study\_log.txt"}. Remember: \textbf{no spaces} around the \texttt{=}.
    \item \textbf{Append a timestamped entry.} Add a line such as \textit{\textless name\textgreater{} studied on \textless date\textgreater} to the log using \texttt{\$1}, command substitution \texttt{\$(date)}, and \texttt{>{}>} (append, \textbf{not} \texttt{>}, which would erase the file each run).
    \item \textbf{Count and report.} Use \texttt{wc -l} to count the lines in the log, then print a message with the total.
\end{enumerate}
\textbf{Optional stretch (not required for full credit):} add an \texttt{if} with the \texttt{-ge} test operator that prints an encouraging message once you've logged 3 or more sessions, and a \texttt{for} loop that lists your existing entries.

% --- SUBMITTING ---
\section{How to Submit}
\begin{enumerate}
    \item Run your script at least twice so \texttt{study\_log.txt} has entries in it.
    \item Zip your \texttt{hw7-studylog} folder, with your finished \texttt{studylog.sh} and the \texttt{study\_log.txt} it produced included:
    \begin{itemize}
        \item \textbf{Windows:} Right-click the \texttt{hw7-studylog} folder \(\rightarrow\) \textbf{Send to} \(\rightarrow\) \textbf{Compressed (zipped) folder}.
        \item \textbf{Mac:} Right-click the \texttt{hw7-studylog} folder \(\rightarrow\) \textbf{Compress ``hw7-studylog''}.
    \end{itemize}
    \item Go to Ed Discussion and open the \textbf{Week 7} submission thread.
    \item Create a new post, attach your zipped folder, and submit.
\end{enumerate}

% --- CHECKLIST ---
\section{Checklist Before You're Done}
\begin{itemize}
    \item \texttt{\#!/bin/bash} is still the first line of the script - 10\%
    \item Running with no name prints a usage message and exits non-zero - 20\%
    \item The log file name is stored in a variable (no spaces around \texttt{=}) - 15\%
    \item A timestamped entry is appended using \texttt{\$1}, \texttt{\$(date)}, and \texttt{>{}>} - 25\%
    \item The total number of sessions is counted with \texttt{wc -l} and printed - 20\%
    \item You ran the script at least twice and \texttt{study\_log.txt} shows your entries - 10\%
\end{itemize}

% --- GRADING ---
\section{Grading}
This assignment is graded on completion of the checklist above. Each item is worth the percentage listed next to it, and the items add up to 100\% of your grade on this assignment.

\end{document}
